\documentclass[11pt]{article}
\usepackage[margin=0.8in]{geometry}
\usepackage[T1]{fontenc}
\usepackage{lmodern}
\usepackage{xcolor}
\usepackage{hyperref}
\usepackage{enumitem}
\usepackage{longtable}
\usepackage{fancyvrb}
\usepackage{titlesec}
\definecolor{ink}{HTML}{302719}
\definecolor{tan}{HTML}{F3EBDD}
\definecolor{blue}{HTML}{175C86}
\hypersetup{colorlinks=true,linkcolor=blue,urlcolor=blue}
\titleformat{\section}{\Large\bfseries\color{ink}}{\thesection}{0.6em}{}
\titleformat{\subsection}{\large\bfseries\color{ink}}{\thesubsection}{0.6em}{}
\newcommand{\step}[1]{\subsection*{Step #1}}
\newcommand{\button}[1]{\fcolorbox{black!15}{tan}{\textsf{#1}}}
\newcommand{\code}[1]{\texttt{\small #1}}
\newenvironment{terminal}{\Verbatim[fontsize=\small,frame=single,framesep=4pt]}{\endVerbatim}
\setlist{itemsep=3pt,topsep=3pt}
\begin{document}
\begin{center}
 {\Huge\bfseries POTW Deployment Guide}\\[6pt]
 {\large A click-by-click guide for Namecheap, Resend, Render, and PostgreSQL}\\[10pt]
 \textit{Prepared for the Purdue Math Club Problem of the Week app}
\end{center}

\begin{center}\colorbox{tan}{\parbox{0.92\linewidth}{\textbf{Security rule:} Never paste a database password, Resend API key, JWT secret, or complete connection URL into GitHub, screenshots, Discord, or chat. Put secrets only in Render environment variables.}}\end{center}

\tableofcontents
\newpage

\section{What you are deploying}
The repository is a single project (a monorepo), but Render will run two services from it:
\begin{itemize}
 \item \textbf{Static Site:} builds the React files in \code{src/} and serves the website.
 \item \textbf{Web Service:} runs the Express API in \code{server/}. It handles login, email, submissions, grading, hints, and database access.
 \item \textbf{PostgreSQL:} stores users, problems, submissions, grades, seasons, and hint requests.
\end{itemize}
The request path is: browser $\rightarrow$ React Static Site $\rightarrow$ Express Web Service $\rightarrow$ PostgreSQL.

\section{Before clicking anything}
Have these ready:
\begin{enumerate}
 \item A GitHub repository containing \code{src/}, \code{server/}, \code{db/}, \code{public/}, and \code{package.json}.
 \item A Namecheap account and the domain \code{purdue-math-potw.com}.
 \item A Resend account.
 \item A Render account connected to the GitHub account.
 \item PostgreSQL command-line tools (\code{psql}) installed locally for one-time migrations.
\end{enumerate}
The repository must not contain \code{.env}, \code{node\_modules/}, or \code{dist/}.

\section{Save and push the project}
Open PowerShell and run:
\begin{terminal}
cd "C:\textbackslash Users\textbackslash YOUR\_NAME\textbackslash Documents\textbackslash GitHub\textbackslash POTW"
git add .
git commit -m "Prepare POTW for deployment"
git push origin main
\end{terminal}
If Git asks for a login, complete the GitHub sign-in. If the push fails, stop and fix that before continuing.

\section{Namecheap and Resend email setup}
\step{Create the Resend domain entry}
\begin{enumerate}
 \item Go to \url{https://resend.com} and sign in.
 \item In the left menu click \button{Domains}.
 \item Click \button{Add Domain}.
 \item Enter \code{purdue-math-potw.com} and click \button{Add}.
 \item Leave this Resend page open. It displays the DNS records you must copy.
\end{enumerate}

\step{Add DNS records in Namecheap}
\begin{enumerate}
 \item Go to \url{https://www.namecheap.com} and sign in.
 \item Click \button{Domain List} in the left sidebar.
 \item Find \code{purdue-math-potw.com} and click \button{Manage}.
 \item Click the \button{Advanced DNS} tab.
 \item Under \button{Host Records}, click \button{Add New Record}.
 \item For every record shown by Resend, choose the same type and paste its exact Host and Value.
 \item Set TTL to \button{Automatic}, then click the green checkmark/save icon.
\end{enumerate}
For a Namecheap Host field, use \code{@} when Resend means the root domain. If Resend gives a full hostname, Namecheap may only need the part before \code{.purdue-math-potw.com}. Do not guess: use Resend's displayed values. Do not create a second SPF record if one already exists.

\step{Verify in Resend}
\begin{enumerate}
 \item Return to Resend \button{Domains}.
 \item Click \code{purdue-math-potw.com}.
 \item Click \button{Verify DNS Records}.
 \item Wait until the domain status says \button{Verified}.
\end{enumerate}

\step{Create the Resend API key}
\begin{enumerate}
 \item In Resend click \button{API Keys}.
 \item Click \button{Create API Key}.
 \item Name it \code{POTW Production}.
 \item Select \button{Sending access} and restrict it to \code{purdue-math-potw.com} when offered.
 \item Click \button{Create Key}.
 \item Copy the complete key beginning with \code{re\_}. Resend only displays it once.
\end{enumerate}

\section{Create Render PostgreSQL}
\begin{enumerate}
 \item Go to \url{https://dashboard.render.com}.
 \item Click \button{New} in the upper-right corner.
 \item Select \button{PostgreSQL}.
 \item Name it \code{potw-db}.
 \item Choose the same region you will use for the backend (Ohio is appropriate for Purdue users).
 \item Choose a paid production plan if you need reliable always-on production behavior.
 \item Click \button{Create Database}.
 \item Wait until the database status says \button{Available}.
 \item Open the database, click \button{Connect}, and copy the \textbf{External Database URL} for local migrations.
\end{enumerate}
Use the database's \textbf{Internal Database URL} later in the Render backend service. Render recommends internal connections for services in the same region.

\section{Unverified signup behavior}
The app intentionally does not let an abandoned, unverified signup reserve a username or email forever. When a new registration matches an existing \textbf{unverified} username or email, the old unverified account is replaced and a fresh verification email is sent. A verified account is never replaced; it correctly receives a ``username or email already registered'' error. This also invalidates the old verification link, so use only the newest email link.

To test this after deployment:
\begin{enumerate}
 \item Register a username and email, but do not click its verification link.
 \item Register again with the same username and email.
 \item Confirm the second registration succeeds and a new email arrives.
 \item Click only the newest link; the old link should be invalid.
 \item Verify the account and then try registering the same username again; this time it must be rejected.
\end{enumerate}

\section{Run the database migrations}
These commands create the schema and apply all features. Run them once, in this exact order.

\step{Open PowerShell in the repository}
\begin{terminal}
cd "C:\textbackslash Users\textbackslash YOUR\_NAME\textbackslash Documents\textbackslash GitHub\textbackslash POTW"
$prodDbUrl = Read-Host "Paste the Render External Database URL"
\end{terminal}
Paste the URL only when PowerShell prompts you. It will not be printed in this guide.

\step{Run each file}
\begin{terminal}
psql $prodDbUrl -f db/001_schema.sql
psql $prodDbUrl -f db/002_sample_problems.sql
psql $prodDbUrl -f db/003_admin_features.sql
psql $prodDbUrl -f db/004_submission_options.sql
psql $prodDbUrl -f db/005_account_security.sql
psql $prodDbUrl -f db/006_problem_release_times.sql
psql $prodDbUrl -f db/007_problem_due_times.sql
psql $prodDbUrl -f db/008_default_due_time.sql
psql $prodDbUrl -f db/009_align_seed_season.sql
psql $prodDbUrl -f db/010_hint_responses.sql
\end{terminal}
\step{Verify}
\begin{terminal}
psql $prodDbUrl -c "\\dt"
psql $prodDbUrl -c "SELECT COUNT(*) AS problems FROM problems; SELECT COUNT(*) AS seasons FROM seasons;"
\end{terminal}
You should see tables including \code{users}, \code{problems}, \code{submissions}, \code{seasons}, \code{problem\_proposals}, and \code{hint\_requests}.\\

\section{Create the Render backend}
\begin{enumerate}
 \item In Render click \button{New} $\rightarrow$ \button{Web Service}.
 \item Select the GitHub repository and click \button{Connect}.
 \item Leave \button{Root Directory} blank. The project is now at the repository root.
 \item Set Runtime to \button{Node}.
 \item Set Build Command to \code{npm install}.
 \item Set Start Command to \code{npm run start:server}.
 \item Click \button{Create Web Service}.
\end{enumerate}

\step{Add backend environment variables}
Open the new service, click \button{Environment} in the sidebar, then click \button{Add Environment Variable}. Add each row:
\begin{longtable}{p{0.28\linewidth}p{0.62\linewidth}}
\code{NODE\_ENV} & \code{production}\\
\code{DATABASE\_URL} & Render PostgreSQL \textbf{Internal Database URL}\\
\code{JWT\_SECRET} & A new random value generated locally\\
\code{FRONTEND\_URL} & Temporarily \code{https://example.com}\\
\code{TRUST\_PROXY} & \code{1}\\
\code{SMTP\_HOST} & \code{smtp.resend.com}\\
\code{SMTP\_PORT} & \code{587}\\
\code{SMTP\_SECURE} & \code{false}\\
\code{SMTP\_USER} & \code{resend}\\
\code{SMTP\_PASS} & Complete Resend API key\\
\code{EMAIL\_FROM} & \code{Purdue Math Club POTW <noreply@purdue-math-potw.com>}\\
\end{longtable}
Generate the JWT value in PowerShell:
\begin{terminal}
[Convert]::ToBase64String([Security.Cryptography.RandomNumberGenerator]::GetBytes(48))
\end{terminal}
Click \button{Save Changes}; Render redeploys the service.

\step{Test the backend}
Copy the backend's Render URL, then open:
\begin{terminal}
https://YOUR-BACKEND.onrender.com/api/health
\end{terminal}
The page must display \code{\{"status":"ok"\}} before you create the frontend.

\section{Create the Render frontend}
\begin{enumerate}
 \item In Render click \button{New} $\rightarrow$ \button{Static Site}.
 \item Select the same GitHub repository and click \button{Connect}.
 \item Leave \button{Root Directory} blank.
 \item Set Build Command to \code{npm install \&\& npm run build}.
 \item Set Publish Directory to \code{dist}.
 \item Under Environment Variables add \code{VITE\_API\_URL}. Set its value to the backend URL:
 \begin{terminal}
 VITE_API_URL=https://YOUR-BACKEND.onrender.com
 \end{terminal}
 \item Click \button{Create Static Site}.
\end{enumerate}

\step{Add the SPA rewrite}
In the frontend service click \button{Redirects/Rewrites} (or \button{Settings} $\rightarrow$ \button{Redirects and Rewrites}), click \button{Add Rule}, and enter:
\begin{longtable}{p{0.25\linewidth}p{0.3\linewidth}p{0.3\linewidth}}
Source & Destination & Action\\
\code{/*} & \code{/index.html} & \code{Rewrite}\\
\end{longtable}
Save the rule and redeploy if Render asks.

\section{Connect the two services}
Copy the frontend URL, for example \code{https://potw-frontend.onrender.com}. Return to the backend service, open \button{Environment}, and replace the temporary value with:
\begin{terminal}
FRONTEND_URL=https://YOUR-FRONTEND.onrender.com
\end{terminal}
Save and redeploy the backend. This allows browser requests from the frontend to reach the API.

\section{First production test}
Run these tests in order:
\begin{enumerate}
 \item Open the frontend Render URL.
 \item Register a test account.
 \item Confirm the verification email arrives.
 \item Log in and submit a solution.
 \item Log in as the admin and open \button{Admin} $\rightarrow$ \button{Grading}.
 \item Grade the submission and confirm the leaderboard updates.
 \item Open \button{Hints}, request a hint, and respond from the admin proposals tab.
 \item Confirm the solver sees the response.
 \item Test file upload, profile history, password reset, and anonymous mode.
 \item Refresh \code{/leaderboard}, \code{/profile}, and \code{/hints}; each must still load because of the rewrite.
\end{enumerate}

\section{Common problems}
\begin{longtable}{p{0.29\linewidth}p{0.62\linewidth}}
No Environment tab & You opened the database or Static Site. Open the backend \textbf{Web Service}.\\
Emails fail with 403 & The domain is not verified, the sender does not use the verified domain, or you used a shortened API key.\\
CORS error & Set backend \code{FRONTEND\_URL} to the exact frontend URL, save, and redeploy.\\
Frontend API errors & Check \code{VITE\_API\_URL}; it must be the backend URL with no trailing \code{/api}.\\
Refresh gives 404 & Add the Static Site rewrite \code{/*} $\rightarrow$ \code{/index.html} as a Rewrite.\\
Database connection fails & Use the Render Internal URL in the backend and confirm the database is Available.\\
Migrations fail with relation exists & Usually safe; many migrations are idempotent. Read the first actual ERROR line.\\
\end{longtable}

\section{After launch}
\begin{itemize}
 \item Enable paid PostgreSQL backups/recovery appropriate for your data.
 \item Rotate the database credential and Resend key if either was exposed.
 \item Add a custom website domain later; Render's generated URL works without one.
 \item Keep the backend service running continuously because its scheduler archives expired problems.
 \item Never put secrets in frontend variables; anything beginning with \code{VITE\_} is visible in the browser build.
\end{itemize}

\section{Useful official documentation}
\begin{itemize}
 \item Render deployment: \url{https://render.com/docs/your-first-deploy}
 \item Render Postgres: \url{https://render.com/docs/postgresql-creating-connecting}
 \item Resend domains: \url{https://resend.com/docs/dashboard/domains/introduction}
 \item Resend API keys: \url{https://resend.com/docs/dashboard/api-keys/introduction}
 \item Resend SMTP: \url{https://resend.com/docs/send-with-smtp}
\end{itemize}
\end{document}
